\section{Introduction}
\label{sec:introduction-and-main-results}

Let \(\rho\) be a non-Gaussian probability law and
\(\mu=\rho^{\otimes\N}\).  We ask which bounded invertible operators on
\(\ell^2\) preserve the measure class of \(\mu\) when their rows act on the
coordinate sequence.  Kakutani's theorem answers the corresponding question
for coordinatewise transformations, but general linear mixing destroys the
independence on which its product criterion relies.  Fixed finite output
marginals do not resolve the problem: they may all have positive densities
even when the full image law is singular.

For finite-variance coordinate laws satisfying the regularity assumptions
below, the answer has the same operator-ideal form as the
Feldman--H\'ajek criterion, but with a discrete non-Gaussian stabilizer.  The
image law is equivalent to \(\mu\) precisely when the mixing operator is a
Hilbert--Schmidt perturbation of an allowed signed coordinate permutation.
We also show that this operator condition is detected exactly by the
Hellinger energies of the actual finite output marginals.

This classification has two sharp boundaries.  Infinite variance changes the
metric cost: for symmetric stable products, off-diagonal columns accumulate
with exponent \(\alpha\), not quadratically.  Zeros of a finite-variance
density create a different obstruction: Hilbert--Schmidt proximity still
controls accumulated Hellinger loss, but equivalence additionally requires
the rows of the operator and its inverse to avoid the zero set.  Thus tail
geometry and support geometry enter the classification in distinct ways.

\subsection{The finite-variance classification}
\label{sec:main-results}

Suppose \(\rho(dx)=f(x)\,dx\), put \(g=\sqrt f\), and assume
\begin{equation}
 \begin{gathered}
 f>0\quad\text{Lebesgue-almost everywhere},\qquad
 g\in W^{1,2}(\R),\qquad xg'\in L^2(\R),\\
 \E X=0,\qquad \E X^2=1,\qquad
 f\ne(2\pi)^{-1/2}e^{-x^2/2}.
 \end{gathered}
 \label{eq:main-assumptions}
\end{equation}
Derivatives are weak derivatives.  These hypotheses impose finite location
and scale Fisher information but no moment above order two; in particular,
cusps are allowed.

Put \(H=\ell^2(\N;\R)\), let \((e_i)_{i\ge1}\) be its coordinate basis, and
write \(\mathcal S_2(H)\) for the Hilbert--Schmidt ideal.  Define
\begin{equation}
 \begin{aligned}
 \mathcal P=\{Q:{}&(Qx)_i=\varepsilon_i x_{\pi(i)},\quad
       \pi\text{ a permutation of }\N,\\
       &\varepsilon_i\in\{-1,1\},\quad
       \Law(\varepsilon_iX)=\rho\text{ for every }i\}.
 \end{aligned}
 \label{eq:allowed-permutations}
\end{equation}
Thus \(\mathcal P\) is the full signed-permutation group when \(\rho\) is
symmetric and the ordinary permutation group otherwise.  For
\(T\in\mathcal B(H)\), let \(\Phi_T\) be the measurable row-series
realization constructed in
\cref{lem:measurable-row-series-realization}, and set
\(\nu_T=(\Phi_T)_\#\mu\).  A sample from \(\mu\) is not assumed to lie in
\(H\).

For probability measures \(\lambda,\nu\), write
\begin{equation}
 \Aff(\lambda,\nu)
 =\int\sqrt{\frac{d\lambda}{dm}\frac{d\nu}{dm}}\,dm,
 \qquad
 h(\lambda,\nu)^2=2-2\Aff(\lambda,\nu),
 \label{eq:affinity-convention}
\end{equation}
where \(m\) dominates both measures.  Let \(\mu_n=\rho^{\otimes n}\), let
\(\pi_n\) extract the first \(n\) coordinates, and define
\begin{equation}
 E_n(T)=-2\log\Aff\bigl(\mu_n,(\pi_n)_\#\nu_T\bigr).
 \label{eq:projected-hellinger-energy}
\end{equation}
These energies use the true output rows of \(T\), rather than principal matrix
corners.

\begin{theorem}[Bounded invertible linear mixing]
\label{thm:measure-equivalence}
Under \eqref{eq:main-assumptions}, every \(T\in\GL(H)\) satisfies
\begin{equation}
 \boxed{\begin{aligned}
 \nu_T\sim\mu
 &\quad\Longleftrightarrow\quad
 \exists Q\in\mathcal P:\ TQ^{-1}-I\in\mathcal S_2(H)\\
 &\quad\Longleftrightarrow\quad \sup_nE_n(T)<\infty.
 \end{aligned}}
 \label{eq:measure-equivalence}
\end{equation}
If these conditions fail, then \(\nu_T\perp\mu\) and
\(E_n(T)\uparrow\infty\).  Moreover, \(\nu_T=\mu\) if and only if
\(T\in\mathcal P\).
\end{theorem}

The theorem identifies both the exact symmetry group and the full
measure-class neighborhood around it.  The projected-energy equivalence is
not a finite-dimensional surrogate: bounded energy first detects a nonzero
absolutely continuous component, and an orbit dichotomy upgrades that overlap
to equivalence.

Two consequences separate the algebraic and analytic parts of the result.

\begin{corollary}[Moment-only stabilizer]
\label{cor:moment-only-stabilizer}
If \(\rho\) is non-Gaussian, centered, and has variance one, then the exact
linear stabilizer of \(\rho^{\otimes\N}\) is \(\mathcal P\).  No density or
Fisher information assumption is needed.
\end{corollary}

\begin{corollary}[Orthogonal transformations]
\label{cor:orthogonal-equivalence}
If \(f>0\) almost everywhere, \(\sqrt f\in W^{1,2}(\R)\), and \(f\) is
non-Gaussian, centered, and has variance one, then every orthogonal \(U\)
satisfies
\[
 \nu_U\sim\mu
 \quad\Longleftrightarrow\quad
 U-Q\in\mathcal S_2(H)\quad\text{for some }Q\in\mathcal P.
\]
Thus scale Fisher information is not needed for orthogonal mixing.
\end{corollary}

\subsection{Relation to earlier criteria}
\label{sec:history}

Kakutani reduced equivalence of infinite products to summability of
one-coordinate Hellinger losses~\cite{kakutani1948equivalence}; Euclidean and
Lie-group perturbations retain this tensor structure and recover Fisher
information as the local metric
\cite{steele1986fisher,marques1992quasi}.  Gaussian linear mixing is governed
instead by the Cameron--Martin and Feldman--H\'ajek theories
\cite{cameron1944transformations,feldman1958equivalence,hajek1958property,
shepp1966radon,ramer1974nonlinear,bogachev1998gaussian}.  There the full
orthogonal group preserves the standard product law.  For a non-Gaussian
i.i.d. product, exact independence reduces this group to \(\mathcal P\), and
\cref{thm:measure-equivalence} determines the entire bounded invertible
measure-class orbit around that discrete stabilizer.  This is distinct from
the existence theory for non-Gaussian quasi-invariant measures on linear
spaces~\cite{feldman1961examples,takahashi1975quasi}.

\subsection{Two boundary regimes}
\label{sec:boundary-regimes}

The finite-second-moment hypothesis determines the quadratic accumulation
law.  For the standard symmetric \(\alpha\)-stable product,
\(0<\alpha<2\), \cref{thm:complete-criterion-symmetric-stable-coordinates}
shows, within the compatible stable row-series class, that equivalence holds
precisely when an allowed signed permutation reduces the operator to a matrix
\(M=(m_{ij})\) satisfying
\[
 \sum_i|m_{ii}-1|^2
 +\sum_j\left(\sum_{i\ne j}|m_{ij}|^2\right)^{\alpha/2}<\infty.
\]
The first term is quadratic, whereas the second is an
\(\ell^\alpha(\ell^2)\) column cost.  Hence an everywhere-positive density
and finite Fisher information do not by themselves make Hilbert--Schmidt
proximity sufficient.

When a finite-variance density may vanish, the same Fourier/Fisher argument
still gives Hilbert--Schmidt necessity.  The exact extension in
\cref{thm:linear-equivalence-group} adds the condition that every row sum of
both \(T\) and \(T^{-1}\) lies almost surely in \(\{f>0\}\).  For bounded or
one-sided coordinate laws, support geometry is stronger still:
\cref{thm:bounded-and-half-line-classification} reduces the equivalence group
to explicit signed-permutation or permutation classes without density or
Fisher assumptions.

\subsection{Proof strategy and organization}
\label{sec:proof-architecture}

The finite-variance proof has three main steps.  First, non-Gaussianity is
encoded by the rotational Fourier defect
\[
 R(\xi,\eta)=\eta\phi'(\xi)\phi(\eta)
              -\xi\phi(\xi)\phi'(\eta),
 \qquad \phi(u)=\E e^{iuX}.
\]
The determinant of the two weighted Fourier jets
\((\xi\phi(\xi),\phi'(\xi))\) and
\((\eta\phi(\eta),\phi'(\eta))\) vanishes at every pair only for a Gaussian
law.  A finite set of frequencies then yields bounded tests for all matrix
directions.  Combined with the matrix Fisher operator, these tests give local
Hellinger coercivity in Frobenius norm with constants independent of the
dimension.

Second, the operator action must be realized on \(\R^{\N}\), since a product
sample need not belong to \(\ell^2\).  A common measurable carrier makes row
series composable and allows inverse maps and likelihood cocycles to be used
without incompatible null sets.  Translation ergodicity gives the
equivalent--singular dichotomy and identifies bounded projected energy with
equivalence.

Third, absolute continuity initially gives only asymptotic matching of remote
rows with allowed signed coordinates.  If the squared matching errors were
not summable, tail-coordinate swaps would package a fixed amount of error into
a small Hilbert--Schmidt perturbation.  Local coercivity detects this
perturbation, while absolute continuity makes the same remote swap
asymptotically invisible.  The contradiction proves Hilbert--Schmidt
necessity.

\Cref{sec:rotational-fourier-defect-and-hellinger-coercivity} develops the
dimension-free Fourier and Hellinger geometry, and
\cref{sec:finite-variance-equivalence-and-global-rigidity} proves the
finite-variance classification.  The stable and zero-set/support regimes are
treated in
\cref{sec:symmetric-stable-products-and-anisotropic-geometry,sec:zero-set-row-condition-and-support-geometry}.
We use \(A^*\) for the adjoint, \(P_n\) for coordinate projection on \(H\),
and \(\|\cdot\|_F\) for the Frobenius norm.
